\section{Canonical Reception Without Retrojection}

Romans 4 makes Abraham the exemplar of faith in the God who gives life to the
dead. Paul's phrase ``against hope believed in hope'' resonates strongly with
an unseen-star meditation, but his argument concerns the barrenness of the
parents and the reliability of the promise. The resonance is legitimate as
canonical reflection; it is not evidence that Paul reconstructed verse 5 as a
daylight scene.

Galatians 3 concentrates on promise, faith, law, and the seed who is Christ.
Hebrews 11 emphasizes obedient departure, sojourning, promised offspring, and
confidence in resurrection. None supplies an hour for the star command.
Together they establish something more durable than the disputed chronology:
the Church receives Genesis 15 within a theology of divine fidelity and
faithful response.

This is why the chronological question is worth asking but cannot bear the
whole doctrine placed upon it. The doctrine does not fall if verse 5 is night,
vision, or an omitted interval. Conversely, the doctrine does not convert a
possible daylight scene into a historical certainty.

